\clearpage
\begin{summarybox}
\textbf{\textcolor{CSummary}{核心口诀与做题流程}}
\begin{center}
\begin{minipage}{\linewidth}\centering
\begin{tikzpicture}[x=1.95cm,y=1.95cm,>=stealth,font=\small,line cap=round,line join=round]
\draw[->,gray!75] (-0.35,0)--(3.96,0) node[right=3pt]{$x$};
\draw[->,gray!75] (0,-0.35)--(0,3.96) node[above=3pt]{$y$};
\node at (-.3,.4) {$O$};
\draw[gray,dashed] (-0.15,-0.15)--(3.8,3.8) node[above left=3pt]{$y=x$};

\draw[gray!55,densely dotted] (1,0)--(1,3)--(0,3) (3,0)--(3,1)--(0,1);
\draw[teal,thick] (1,3)--(3,1);
\fill[teal] (1,3) circle(2pt) node[above left=5pt]{$P(1,3)$};
\fill[blue] (3,1) circle(2pt) node[below right=5pt]{$P'(3,1)$};
\fill (2,2) circle(1.8pt) node[right=12pt]{$M(2,2)$};
\draw (1.84,2.16)--(1.68,2)--(1.84,1.84);

\end{tikzpicture}\par\smallskip
{\small $PP^\prime\perp y=x,\qquad PM=MP^\prime$}\par
{\small 中点在镜线上，横纵坐标互换；两轴采用相同单位长度。}\par\smallskip
\end{minipage}
\end{center}


\textbf{口诀：}关于 $y=x$ 对称，横纵互换；函数可逆，互为反函数；自己反自己，复合两次回原值。
\[
 (a,b)\leftrightarrow(b,a),\quad F(x,y)=0\leftrightarrow F(y,x)=0.
\]
\begin{enumerate}[leftmargin=*,itemsep=.4em]
\item 先识别对称轴是否为 $y=x$，然后交换横纵坐标与全部限制条件。
\item 若要求反函数，先判断一一对应，再解出新纵坐标，标明定义域和值域。
\item 若要求自身对称，检查 $f(D)\subseteq D$，证明全域 $f(f(x))=x$。
\item 若要求交点，先看严格递增条件；否则使用公共点的复合条件。
\end{enumerate}
\end{summarybox}
\clearpage
\begin{summarybox}
\textbf{\textcolor{CSummary}{最常见模型速查}}
\begin{center}
\begin{minipage}{\linewidth}\centering
\begin{tikzpicture}[x=1.7cm,y=1.7cm,>=stealth,font=\small,line cap=round,line join=round]
\draw[->,gray!75] (-0.35,0)--(4.76,0) node[right=3pt]{$x$};
\draw[->,gray!75] (0,-0.35)--(0,4.76) node[above=3pt]{$y$};
\node at (-.3,.4) {$O$};
\draw[gray,dashed] (-0.15,-0.15)--(4.6,4.6) node[above left=3pt]{$y=x$};

\draw[teal,very thick,domain=0:2,samples=100] plot(\x,{\x*\x});
\draw[blue,very thick,domain=0:2,samples=100] plot({\x*\x},\x);
\draw[orange!80!black,densely dotted] (2,4)--(4,2);
\fill[teal] (2,4) circle(2pt) node[above left=5pt]{$(2,4)$};
\fill[blue] (4,2) circle(2pt) node[below right=5pt]{$(4,2)$};

\end{tikzpicture}\par\smallskip
{\small \textcolor{teal}{\rule[.5ex]{1em}{.8pt}\ $y=x^2,\ x\ge0$}\qquad\textcolor{blue}{\rule[.5ex]{1em}{.8pt}\ $y=\sqrt{x},\ x\ge0$}}\par
{\small 选取非负分支后，每个输出恰好对应一个输入。}\par\smallskip
\end{minipage}
\end{center}


\begin{center}\renewcommand{\arraystretch}{1.4}
\begin{tabular}{c|c}
原函数（或曲线） & 镜像\\\hline
$a^x$（$a>0,a\ne1$） & $\log_a x$\\
$x^2$（$x\ge0$） & $\sqrt{x}$\\
$x^3$ & $\sqrt[3]{x}$\\
$kx+b$（$k\ne0$） & $(x-b)/k$\\
$k/x$（$k\ne0$） & 自身\\
$b-x$ & 自身\\
$y^2=2px$ & $x^2=2py$
\end{tabular}\end{center}
\textbf{最后提醒：}反函数与原函数定义域、值域互换；严格单调方向保持；斜率取倒数而不是负倒数；导数必须取对应点且原导数非零。严格递增时，交点可无遗漏地通过 $f(x)=x$ 求得；其他情形需另外判断。
\end{summarybox}
